\begin{abstract}
  Current time-series foundation models forecast in value space,
  following the language-modelling recipe of next-token prediction.
  Because the observed series is one realisation of a stochastic
  process, a value-space loss asks the model to reproduce intrinsic
  noise that no model can predict. We propose \emph{Contrastive
  Forecasting}, a training framework in which an encoder and a causal
  forecaster are trained jointly to predict the next patch in latent
  space through contrastive learning, with negatives drawn along the
  cross-temporal, cross-channel and cross-batch axes. The framework
  is decoupled from the choice of encoder and forecaster
  architecture. On synthetic stationary ARMA processes, the latents
  produced by the backbone carry enough information about the
  generating process for a small head to recover its parameters.
\end{abstract}
